\chapter{Economic Data with Python}\label{ch:b1:economic-data-python}

Three analysts are asked how fast the euro-area economy grew from 2000 to
2023. The first divides GDP in current euros in 2023 by GDP in 2000 and finds
$3.24\%$ a year. The second uses GDP in chained volumes and finds $1.23\%$. The
third divides by population and finds $0.95\%$. % computed: code/bachelor-1/28-economic-data-python/outputs/per_person_growth.txt
All three are right; they answer three different questions, and a report
that prints one of them without saying which has made an error that no
reader can find. This chapter is about the steps between a published
table and a number in a report: loading, joining, converting, adjusting
and writing --- each done by a script that can be run again.

\begin{openingchart}{US consumer prices, month-on-month change (per cent),
2015--2024: not seasonally adjusted (blue) and seasonally adjusted by the
statistical office (red).}{Bureau of Labor Statistics, CPI-U all items
(CUUR0000SA0, CUSR0000SA0)}
\begin{tikzpicture}
\begin{axis}[omtime, xlabel={Year}, ylabel={Change on previous month (\%)}, xmin=2015, xmax=2025, ymin=-1,
  ymax=1.5, width=0.75\linewidth, height=0.38\linewidth]
  \addplot[omDef, thin] table[x=t, y=nsa, col sep=comma] {figdata/bachelor-1/28-economic-data-python/out/cpi_mom.csv};
  \addplot[omThm, thick] table[x=t, y=sa, col sep=comma] {figdata/bachelor-1/28-economic-data-python/out/cpi_mom.csv};
\end{axis}
\end{tikzpicture}
\end{openingchart}

The unadjusted series swings with the calendar: prices rise faster in the
first months of the year and slower at its end. Over the 120 months its
standard deviation is $0.345$ points, against $0.287$ for the adjusted series.
% figdata: cpi_stats
Read without adjustment, a fall in December looks like news; it is the
season.

\begin{recall}
Python, lists, dictionaries, functions, reading a file and first tables
were met in the computer science series. Growth rates and logarithms are
\cref{ch:b1:models-and-data}; price indices, deflators and chained
indices are \cref{ch:b1:consumer-welfare}.
\end{recall}

\section{Three shapes of data}\label{sec:b1:economic-data-python:shapes}

\begin{definition}[Cross-section, time series, panel data]\label{def:b1:economic-data-python:shapes}
\index{cross-section}\index{time series}\index{panel data}
A \emph{cross-section} observes many units at one date; a \emph{time series}
observes one unit at many dates; \emph{panel data} observe many units at many
dates.
\end{definition}

\begin{definition}[Unit of observation]\label{def:b1:economic-data-python:unit}
\index{unit of observation}
The \emph{unit of observation} is what one row of a table describes: a
country in a year, a household, a firm in a month.
\end{definition}

\begin{definition}[Merge key]\label{def:b1:economic-data-python:key}
\index{merge key}
A \emph{merge key} is the set of columns that identifies a \omterm{def:b1:economic-data-python:unit}{unit of
observation}, such as \texttt{(iso3, year)}: no two rows of a table share it,
and two tables are joined by matching it.
\end{definition}

A panel can be stored \emph{long}, one row per unit and date, or \emph{wide},
one row per unit and one column per date. Long tables are what programs
join and filter; wide tables are what people read.

\omcode{code/bachelor-1/28-economic-data-python/listings/first_look.py}{1}{8}{A first look at a
table: its shape, its column types and its first rows.}\label{lst:b1:economic-data-python:first-look}

\omoutput{code/bachelor-1/28-economic-data-python/outputs/first_look.txt}{Eighteen rows: two
countries over nine years. The key is \texttt{(iso3, year)}.}\label{out:b1:economic-data-python:first-look}

Three things are worth reading before any computation. The shape says how
many rows to expect from every later step. The types say that
\texttt{iso3} and \texttt{country} are text and \texttt{year} an integer, so a year
can be compared and sorted. The numbers at the left, 3469 to 3472, are
positions in the full table, not years. The output does not give the units:
\texttt{rgdpe} is in millions of 2021 dollars and \texttt{pop} and \texttt{emp} in
millions of people. They come from the source's documentation, and any script
that divides one column by another must use them.

\omcode{code/bachelor-1/28-economic-data-python/listings/long_and_wide.py}{1}{16}{From long to wide
and back: \texttt{pivot} and \texttt{melt}.}\label{lst:b1:economic-data-python:long-wide}

\omoutput{code/bachelor-1/28-economic-data-python/outputs/long_and_wide.txt}{Twenty-five rows long,
five by five wide, twenty-five again after melting.}\label{out:b1:economic-data-python:long-wide}

\texttt{pivot} takes the key apart: \texttt{iso3} becomes the rows and
\texttt{year} the columns, so each pair may appear only once, or pandas raises
an error. \texttt{melt} puts the key back together, with the former column
names in a column of their own. The round trip returns the 25 rows it
started from, a cheap check that nothing was lost. The wide table is the
one to read. Germany's GDP per person in 2023, 59,854 dollars, was still
below its 2019 level of 60,283; seen from the long table, that comparison
needs a filter.

\section{Loading and joining}\label{sec:b1:economic-data-python:joining}

Two sources rarely cover the same units. A join keeps what matches, and
whatever does not match is lost without a message unless one asks for it.
The argument \texttt{indicator=True} adds a column that says, for each row,
whether its key was found in both tables or in one only.

\begin{omfigure}
\begin{tikzpicture}[font=\footnotesize]
  \draw[thick, omDef] (0,0) circle (1.3);
  \draw[thick, omThm] (1.6,0) circle (1.3);
  \node at (-0.75,0) {\shortstack{World Bank\\only: 69}};
  \node at (0.8,0) {\shortstack{both\\178}};
  \node at (2.2,0) {\shortstack{Penn\\only: 7}};
  \node[anchor=north] at (0.8,-1.4) {key: \texttt{(iso3, year)}, 2019};
\end{tikzpicture}

{\small Joining the World Bank and Penn tables for 2019. The 69 keys found only
in the World Bank table are mostly aggregates (regions, income groups);
the 7 found only in the Penn table are economies the World Bank series
does not cover that year.}
\end{omfigure}

\omcode{code/bachelor-1/28-economic-data-python/listings/join_sources.py}{1}{13}{An outer join with an
indicator: count what did not match.}\label{lst:b1:economic-data-python:join}

\omoutput{code/bachelor-1/28-economic-data-python/outputs/join_sources.txt}{178 matches, 69 keys
only in the World Bank table, 7 only in the Penn table.}\label{out:b1:economic-data-python:join}

The first keys found only in the World Bank table show what most of them
are. AFE and AFW are the two halves of sub-Saharan Africa, ARB is the Arab
world and CEB is central Europe and the Baltics: aggregates, not countries.
AFG and AND are countries that the Penn table does not cover. The seven keys
found only in the Penn table --- Anguilla, Montserrat, the British Virgin
Islands, Taiwan, South Sudan, Venezuela and Yemen --- have no value in this
World Bank series. An inner join would have dropped all 76 keys without a
word.

The commonest silent error is the opposite one: a key that is too short.
Merging two panels on the country alone, forgetting the year, pairs every
year of one table with every year of the other.

\omcode{code/bachelor-1/28-economic-data-python/listings/join_trap.py}{1}{13}{A merge on too few keys,
and the check that would have caught it.}\label{lst:b1:economic-data-python:trap}

\omoutput{code/bachelor-1/28-economic-data-python/outputs/join_trap.txt}{Five years each side:
the short key multiplies the rows by five.}\label{out:b1:economic-data-python:trap}

The arithmetic shows what happened. Over the five years the World Bank
table has 199 economies and the Penn table 185, and 178 are in both. On the
right key the merge keeps the 890 country-years found in both. On the
country alone, each of those country-years is paired with all five years of
the other table: $890 \times 5 = 4{,}450$ rows, every one of which looks like
valid data.

The argument \texttt{validate="one\_to\_one"} makes pandas raise an error if
a key appears twice on either side. It costs one line and turns a silent
error into a loud one.

\section{Units and frequencies}\label{sec:b1:economic-data-python:units}

A series is a level (GDP in euros), a ratio (GDP per person, a debt ratio)
or a rate (inflation). Turning monthly data into yearly data depends on
what the series measures: a flow (spending in a month) is summed; a price
or a rate is averaged; a stock (debt at the end of a month) takes the last
value.

\omcode{code/bachelor-1/28-economic-data-python/listings/monthly_to_yearly.py}{1}{14}{The yearly average
of a monthly price index, the December value, and the published annual
average.}\label{lst:b1:economic-data-python:yearly}

\omoutput{code/bachelor-1/28-economic-data-python/outputs/monthly_to_yearly.txt}{The computed averages
equal the published ones to three decimals.}\label{out:b1:economic-data-python:yearly}

The published annual figure is the average of the twelve months, and the
listing finds it again to the third decimal. In these years of rising
prices the December value is above the average every year, by 1.7 points in
2020 and 4.1 in 2022. A report that compares one year's December with
another year's average mixes two measures.

The helper \texttt{monthly} that these listings import drops the row the
statistical office uses for the annual average (period \texttt{M13}) and
indexes the rest by \texttt{(year, month)}.

\omcode{code/bachelor-1/28-economic-data-python/listings/cpi_tools.py}{1}{19}{A monthly series indexed by
year and month.}\label{lst:b1:economic-data-python:cpi-tools}

This file is a module, not a script: it prints nothing, and the other
listings import its function, so that every one of them filters the period
codes in the same way. The warning it silences concerns a column of the
loader's file that no listing reads.

\begin{definition}[Annualised rate]\label{def:b1:economic-data-python:annualised}
\index{annualised rate}
The \emph{annualised rate} of a change over a period shorter than a year is
the change it would give over a year if repeated: for a monthly change $g$,
$(1 + g)^{12} - 1$; for a quarterly change, $(1 + g)^4 - 1$.
\end{definition}

\omcode{code/bachelor-1/28-economic-data-python/listings/changes.py}{1}{13}{Three ways to report one
month's inflation.}\label{lst:b1:economic-data-python:changes}

\omoutput{code/bachelor-1/28-economic-data-python/outputs/changes.txt}{June 2022: a monthly change
of 1.26\% annualises to 16\%; the twelve-month change is 9\%.}\label{out:b1:economic-data-python:changes}

The same month gives a rise of $1.26\%$, $8.98\%$ or $16.16\%$, depending on
the question. % computed: code/bachelor-1/28-economic-data-python/outputs/changes.txt
An annualised monthly rate is noisy and amplifies the season; a
twelve-month change is smooth but six months late on average.

\section{Seasonal adjustment}\label{sec:b1:economic-data-python:seasonal}

\begin{definition}[Seasonal adjustment, seasonal factor]\label{def:b1:economic-data-python:seasonal}
\index{seasonal adjustment}\index{seasonal factor}
A \emph{seasonal factor} is the typical ratio of a series in a given month
(or quarter) to its trend. \emph{Seasonal adjustment} divides each value by
its month's factor, removing the movements that recur at the same time
every year.
\end{definition}

The simplest method estimates the trend by a centred twelve-month moving
average, with weights $1/24$ on the two ends and $1/12$ on the eleven months
between them, so that every calendar month counts equally. The ratio of
the series to that trend, averaged month by month and scaled to average 1,
gives the factors.

\omcode{code/bachelor-1/28-economic-data-python/listings/seasonal_by_hand.py}{1}{27}{The ratio to a
centred moving average, as a function.}\label{lst:b1:economic-data-python:seasonal}

\omoutput{code/bachelor-1/28-economic-data-python/outputs/seasonal_by_hand.txt}{Prices sit $0.43\%$ above
trend in June and $0.63\%$ below in December; the by-hand series moves with the
official one.}\label{out:b1:economic-data-python:seasonal}

The factors run from 0.9937 in December to 1.0043 in June, a swing of about
one per cent within the year. That is as large as a whole month of the
inflation of 2022, which is why the unadjusted series of the opening chart
swings. In the function, a twelve-month mean averaged with itself shifted
by one month gives the centred trend, with weights $1/24$ at the two ends.
The first and last six months have no centred trend and are dropped. The
official method fills them with forecasts (measure box below).

\begin{omchart}{US consumer prices, month-on-month change (per cent),
2015--2024: seasonally adjusted by the statistical office (red) and by the
listing (black, dashed).}{Bureau of Labor Statistics, CPI-U; adjustment by
\cref{lst:b1:economic-data-python:seasonal}}
\begin{tikzpicture}
\begin{axis}[omtime, xlabel={Year}, ylabel={Change on previous month (\%)}, xmin=2015, xmax=2025, ymin=-1,
  ymax=1.5, width=0.75\linewidth, height=0.36\linewidth]
  \addplot[omThm, thick] table[x=t, y=sa, col sep=comma] {figdata/bachelor-1/28-economic-data-python/out/cpi_mom.csv};
  \addplot[black, thin, dashed] table[x=t, y=hand, col sep=comma] {figdata/bachelor-1/28-economic-data-python/out/cpi_mom.csv};
\end{axis}
\end{tikzpicture}
\end{omchart}

The two adjusted series have a correlation of $0.975$ over the 120 months.
% figdata: cpi_stats
The chart's table is written by a script that imports the listing's
function, so the chart and the listing cannot drift apart.

\begin{measure}[Seasonal adjustment]
The statistical office computes its factors with the X-13ARIMA-SEATS
method, which extends the series with forecasts before averaging, so that
the most recent months have a centred trend, and treats outliers and
calendar effects. Its factors are updated each February, and the
previous five years of adjusted data are revised. % ledger: b1:bls-x13
An adjusted figure is therefore never final: last year's December will
change next February.
\end{measure}

\section{Deflating, indexing, logarithms and growth}\label{sec:b1:economic-data-python:deflating}

A real series divides a nominal one by a price index and multiplies by its
base-year value. An index rebased to year $b$ divides by its value in $b$ and
multiplies by 100. The difference of logarithms is a growth rate in log
points; their mean over a span is the average log growth, close to the
compound rate when growth is small.

\omcode{code/bachelor-1/28-economic-data-python/listings/deflate_and_index.py}{1}{14}{Brenland's real GDP,
an index with 2020 = 100, and its mean log growth.}\label{lst:b1:economic-data-python:deflate}

\omoutput{code/bachelor-1/28-economic-data-python/outputs/deflate_and_index.txt}{Real GDP grew by
$14.3\%$ from 2020 to 2025; the recomputed series equals the published one.}\label{out:b1:economic-data-python:deflate}

The last line is a check: recomputing a published series and finding no
difference is how one learns that the deflator and the base year were the
ones assumed. A deflator chained year by year (\cref{ch:b1:consumer-welfare})
makes real components that do not add up to real GDP: a check that adds up
real consumption, investment, government purchases and net exports then
fails, for a reason that must be stated. The check of the listing cannot
fail that way, because real GDP is recomputed from its own deflator.

\omcode{code/bachelor-1/28-economic-data-python/listings/per_person_growth.py}{1}{19}{The three
growth rates of the chapter's opening.}\label{lst:b1:economic-data-python:per-person}

\omoutput{code/bachelor-1/28-economic-data-python/outputs/per_person_growth.txt}{Prices, then
population, take two points and then a quarter of a point off the
first number.}\label{out:b1:economic-data-python:per-person}

Each line divides by one more thing. Prices take 2.01 points a year off
nominal growth, from 3.24 to 1.23, and population growth another 0.28 points,
leaving 0.95. The function \texttt{yearly} computes a compound rate from the
two end years only, so a recession inside the span does not show in it.
That is a choice of method, and it belongs next to the number when the
number is printed.

\section{A chart is a table a script wrote}\label{sec:b1:economic-data-python:reproducible}

\begin{definition}[Reproducible computation]\label{def:b1:economic-data-python:reproducible}
\index{reproducible computation}
A computation is \emph{reproducible} when running its scripts again on the
same inputs gives the same outputs, byte for byte, with no step done by
hand.
\end{definition}

Every chart of this book is drawn from a CSV file that a tested script
wrote; no number in a chart was typed. When a source is revised, the
scripts run again and every chart and every quoted number follows.

\omcode{code/bachelor-1/28-economic-data-python/listings/write_chart_table.py}{2}{15}{Writing a chart's
table twice and comparing the bytes.}\label{lst:b1:economic-data-python:write}

\omoutput{code/bachelor-1/28-economic-data-python/outputs/write_chart_table.txt}{The same digest twice:
the file is the same.}\label{out:b1:economic-data-python:write}

A digest is a fingerprint of a file's bytes: a change of one character
anywhere gives a different one. Two runs that print the same sixteen
hexadecimal digits wrote the same file. The fixed line ending matters
here. Left to its default, the operating system chooses the end-of-line
character, and two machines could write the same numbers in different
bytes.

\begin{record}[A spreadsheet and two choices]
A widely cited study of public debt and growth (Reinhart and Rogoff 2010)
in 20 advanced economies
over 1946--2009 reported that country-years with public debt above $90\%$ of
GDP had average growth of $-0.1\%$. Working from the authors' own
spreadsheet, a replication found that a coding error left five countries
out of an average, that some available years were excluded, and that the
averages were weighted unusually; corrected, the growth rate was $2.2\%$
(Herndon, Ash and Pollin 2014). % ledger: b1:herndon2014
A formula typed into a cell is invisible on the page that cites its
result. A script that can be rerun, and a check that recomputes a
published number, are how such errors are found.
\end{record}

\section*{Exercises}\label{sec:b1:economic-data-python:exercises}

\begin{exercise}[$\star$]\label{exo:b1:economic-data-python:shapes-ex}
Give the shape, the \omterm{def:b1:economic-data-python:unit}{unit of observation} and the key of the tables printed
by \cref{out:b1:economic-data-python:first-look,out:b1:economic-data-python:long-wide},
and of the monthly series of \cref{lst:b1:economic-data-python:cpi-tools}.
\end{exercise}

\begin{exercise}[$\star$]\label{exo:b1:economic-data-python:yoy}
From \cref{out:b1:economic-data-python:changes}, give the twelve-month change of the
unadjusted and the adjusted index in June 2022. Why do they differ?
\end{exercise}

\begin{exercise}[$\star$]\label{exo:b1:economic-data-python:annualise}
Prices rise by $0.4\%$ in a month. Give the \omterm{def:b1:economic-data-python:annualised}{annualised rate}.
\end{exercise}

\begin{exercise}[$\star$]\label{exo:b1:economic-data-python:aggregate}
Which aggregation from months to years --- mean, sum or last value --- suits
the unemployment rate, exports, the consumer price index, the population
and government debt?
\end{exercise}

\begin{exercise}[$\star\star$]\label{exo:b1:economic-data-python:maddison}
Merge the Penn World Table and the Maddison Project table for 2019 on
\texttt{(iso3, year)}. How many keys match, and how many are found in one
table only? Name two kinds of unit in each group.
\end{exercise}

\begin{exercise}[$\star\star$]\label{exo:b1:economic-data-python:rebase}
Rebase Brenland's consumer price index to 2025 = 100 and compute its average
log growth from 2016 to 2025.
\end{exercise}

\begin{exercise}[$\star\star$]\label{exo:b1:economic-data-python:toy}
A quarterly series is $x_t = 100 + t + s_t$ for $t = 0, \dots, 11$, with a season $s$ of
$+2$, $-1$, $-2$ and $+1$ in quarters 1 to 4. Using the differences from a centred
four-quarter moving average, recover the season.
\end{exercise}

\begin{exercise}[$\star\star$]\label{exo:b1:economic-data-python:falling}
A country's population falls. Predict the sign of the difference between
the growth of its GDP per person and the growth of its GDP.
\end{exercise}

\begin{exercise}[$\star\star$]\label{exo:b1:economic-data-python:hook}
Recompute the three growth rates of \cref{lst:b1:economic-data-python:per-person} for
2000--2019. Which of the three changes most when the pandemic years are
left out?
\end{exercise}

\begin{exercise}[$\star\star\star$]\label{exo:b1:economic-data-python:quarterly}
Write \texttt{adjust\_q}, the quarterly version of \texttt{adjust}, and test it on a
series with a known multiplicative season of $0.98$, $1.01$, $1.03$ and $0.98$ and a
trend of $0.5\%$ a quarter.
\end{exercise}

\begin{exercise}[$\star\star\star$]\label{exo:b1:economic-data-python:silent}
A colleague's export of the Penn table writes country codes in lower
case. A left join of the World Bank table with it runs without error. What
does it return, and which check would catch the error?
\end{exercise}

\begin{exercise}[$\star\star\star$]\label{exo:b1:economic-data-python:base}
A real GDP series was computed with a deflator whose base year is not the
one stated. What simple check would reveal it?
\end{exercise}

\section{Problem: Brenland's Open-Data Portal}\label{sec:b1:economic-data-python:problem}

\begin{problem}[{Weekend problem --- checking a portal's numbers}]\label{pb:b1:economic-data-python:portal}
Brenland's statistical office (invented) publishes its accounts, its trade
with Kestria and its consumer basket, and a headline: ``\omterm{prop:b1:models-and-data:real-growth}{real growth},
2015--2025: $4.7\%$ a year''. Its trade file has one year mistyped. Its monthly
consumer price index for 2015--2024 is built from the yearly index, with a
\omterm{def:b1:economic-data-python:seasonal}{seasonal factor} of $1 + 0.004\sin(2\pi(m - 3)/12)$ in month $m$ and a little noise. Each
answer comes from a script.

\medskip\noindent\textbf{Part I --- Loading and checking.}
\begin{enumerate}
  \item How many rows does the accounts table have, and which years?
  \item Check that $GDP = C + I + G + X - M$ in every year.
  \item Check that the deficit equals spending minus revenue.
  \item How many rows do the trade and basket tables have, and why 130?
  \item What is the price index in 2015, and why check it?
  \item Why run these checks before any analysis?
\end{enumerate}

\medskip\noindent\textbf{Part II --- Joining.}
\begin{enumerate}[resume]
  \item Join the accounts and the portal's trade file on \texttt{year} with an
    indicator. Which year of the accounts finds no match?
  \item Which year appears twice in the trade file?
  \item What would \texttt{validate="one\_to\_one"} have done?
  \item Fix the file, and join again.
  \item What share of Brenland's exports went to Kestria in 2025?
  \item What were those exports worth in kests?
\end{enumerate}

\medskip\noindent\textbf{Part III --- The monthly index.}
\begin{enumerate}[resume]
  \item What is the true \omterm{def:b1:economic-data-python:seasonal}{seasonal factor} of June?
  \item What factor does \texttt{adjust} recover for June?
  \item And for December?
  \item What is the largest error of the recovered factors?
  \item Compare the twelve-month change in December 2024 with the change in
    the yearly average from 2023 to 2024.
  \item Which of the two would you report for 2024, and why?
\end{enumerate}

\medskip\noindent\textbf{Part IV --- The headline.}
\begin{enumerate}[resume]
  \item How did the portal obtain $4.7\%$?
  \item Compute \omterm{prop:b1:models-and-data:real-growth}{real growth} with the deflator.
  \item Compute the mean log growth, and say why it is a little lower.
  \item Compute growth of real GDP per person.
  \item Compute growth of nominal GDP divided by the consumer price index,
    and say why it differs from the answer to question 20.
  \item State the result: the corrected \omterm{prop:b1:models-and-data:real-growth}{real growth} of Brenland, 2015--2025.
\end{enumerate}
\end{problem}
